% ============================================================================================
\section{S3: Quiver Plots and Streamlines (Owner: Parthsarati Samanta, BT2024083)}
\label{sec:s3}
\textbf{Method.} Vectors are the 10\,m wind components from Eq.~\eqref{eq:wind}, i.e.\ the direction the air moves. They are
given as geographic east/north components and passed to Cartopy with \texttt{transform=PlateCarree}, which rotates them
into the Lambert map frame. Arrows use \texttt{pivot='middle'}, and a quiver key gives the 10\,m/s reference length. We did
both assignment options: (a) grid sampling with proportional vs.\ same-length arrows, and (b) streamlines vs.\ quiver with
three seeding strategies. We also added regional zooms of all four landfalls. D7 (Helene landfall) is the main test case
because it has the strongest and most structured winds of the season. Script: \fpath{scivis/s3_quiver_streamlines/run.py}.

\subsection{Grid sampling: stride and aliasing}
The native grid has 483,879 land vectors, about 4\,km apart, far too many to draw. Fig.~\ref{fig:18} shows plain striding
(every $k$-th cell) for $k$ = 10, 20, 40, giving 4,862, 1,219 and 295 arrows. Fig.~\ref{fig:19} compares striding with
$k\times k$ \emph{block averaging}:
\begin{equation}
\bar{\mathbf{v}}_B = \frac{1}{|B\cap\Omega|}\sum_{(i,j)\in B\cap\Omega}(u_{ij},v_{ij}),
\end{equation}
where $\Omega$ is the land mask. Blocks with less than 50\,\% land are dropped. Table~\ref{tab:alias} quantifies the
difference over all eight dates (code in Appendix~\ref{app:code}).

\figfull{Fig_18_quiver_stride.png}{Plain striding on D7: every 10th, 20th and 40th cell (about 40, 80 and 160\,km spacing);
magnitude-proportional arrows over the speed field (GnBu, 0--14\,m/s). Reference arrow 10\,m/s.}{fig:18}

\figfull{Fig_19_quiver_stride_vs_block.png}{Striding vs.\ block averaging in the Helene zoom on D7 for $k=20$ (top) and $k=40$
(bottom). (c) Direction difference between the point sample and the block mean, pooled over all 8 dates, vs.\ spacing.
(f) Speed difference and the ratio $|\bar{\mathbf{v}}|/\overline{|\mathbf{v}|}$, which measures vector cancellation inside a
block.}{fig:19}
\input{generated/tab_s3_alias}

\textbf{Inferences.}
(1) $k=10$ (4,862 arrows) produces a solid black texture. Arrows overlap in the 8--12\,m/s Florida flow, and the speed
background is hidden. $k=40$ (295 arrows) is clean but represents a hurricane-scale circulation over the Southeast with
only a few dozen arrows. $k=20$ (about 80\,km, 1,219 arrows) is the best compromise at CONUS scale. Clutter rises roughly
with $k^{-2}$, while the information that remains depends on the scale of the feature.
(2) Striding \emph{aliases}: one 4\,km cell stands for a 160\,km block. The median direction error against the block mean
grows from 1.1$^\circ$ ($k=5$) to 5.0$^\circ$ ($k=40$), and the 90th percentile from 5$^\circ$ to 21$^\circ$. At $k=40$,
2.5\,\% of sampled arrows point more than 45$^\circ$ away from the local mean flow, mostly in mountains and coastal
convergence zones. Their direction depends on which cell the stride happens to hit.
(3) Block averaging is a low-pass filter, so it is the correct choice before subsampling. In Fig.~\ref{fig:19} the
block-mean arrows over Georgia and the Carolinas turn smoothly with Helene's circulation. The strided arrows at $k=40$
include some arrows across the Appalachians that deviate visibly from their neighbours.
(4) Averaging has a cost: vectors cancel inside a block. The ratio of the block-mean vector's magnitude to the mean
speed falls from 0.999 to 0.975 at $k=40$. Averaged arrows are therefore slightly too short where the flow turns,
for example near a circulation centre. The speed background (from $v_s$, not $|\bar{\mathbf{v}}|$) is kept so this bias
is visible.

\subsection{Magnitude-proportional vs.\ uniform-length arrows}
\figfull{Fig_20_quiver_proportional_vs_uniform.png}{D7 wind with 20$\times$20 (CONUS) and 8$\times$8 (Helene zoom) block means.
(a, c) Arrow length proportional to speed; (b, d) unit-length arrows coloured by speed (plasma, 0--12\,m/s).}{fig:20}

\textbf{Inferences.}
(1) Proportional arrows (Fig.~\ref{fig:20}a, c) make the speed gradient obvious: the Florida peninsula arrows are 3--4 times
longer than those over the Great Basin. But the long arrows overlap their neighbours (the 8$\times$8 Florida zoom becomes a
black mass), and the direction of weak flow (below 2\,m/s across Nevada, Utah and Arizona) is nearly unreadable because
those arrows shrink to dots.
(2) Unit-length arrows (Fig.~\ref{fig:20}b, d) show direction everywhere at equal visual weight. They reveal flow
patterns that the proportional version hides: diffuse, variable flow over the interior West; a continuous belt of
southerly (northward-moving) flow from Texas through the Plains to the Dakotas; and the cyclonic turning of Helene's flow
from south-south-east over Florida to east and south-east over Georgia and the Carolinas.
(3) Colour carries the magnitude in the uniform version. It reads less precisely than length but never overlaps. The
peninsula mean is $(\bar u,\bar v)=(-5.3,+9.5)$\,m/s with mean speed 11.2\,m/s (yellow). The same arrows in panel (c) are
the longest of the map.
(4) Recommendation: use proportional arrows for quantitative comparison at coarse density; use uniform arrows plus colour
to show circulation patterns at high density.

\subsection{Regional zooms of the landfalls}
\figfull{Fig_21_hurricane_zooms.png}{The four landfall days: speed (colour) with 7$\times$7 block-mean vectors. The red X
marks the maximum of smoothed cyclonic relative vorticity $\zeta=\partial v/\partial x-\partial u/\partial y$ over land;
the star marks the maximum daily-mean speed.}{fig:21}
\input{generated/tab_s3_regional}

\textbf{Inferences.}
(1) \textbf{Beryl, D1.} A closed counter-clockwise circulation is centred over the upper Texas coast; the vorticity
maximum, $24\times10^{-5}$\,s$^{-1}$, is at 29.78\,$^\circ$N, 96.06\,$^\circ$W. The strongest daily-mean wind (12.0\,m/s) is
just east of it at 29.44\,$^\circ$N, 95.43\,$^\circ$W (south of Houston, near Freeport), consistent with the right-front quadrant of a
northward-moving storm. Beryl made landfall near Matagorda (about 28.7\,$^\circ$N, 96.0\,$^\circ$W) as a Category~1
hurricane at about 09~UTC on 8~July~\cite{nhc_tcr}. The daily-mean vortex sits about 120\,km further north because the
storm moved inland during the gridMET day, so a daily field places a moving storm along its track, not at landfall.
A broad band of 6--9\,m/s southerly inflow covers Louisiana.
(2) \textbf{Debby, D3.} The circulation centre (vorticity maximum $45\times10^{-5}$\,s$^{-1}$, the strongest of the four) is
over the Florida Big Bend at 30.23\,$^\circ$N, 83.35\,$^\circ$W. That is about 60\,km north of Debby's landfall near Steinhatchee
(Category~1, about 11~UTC on 5~August~\cite{nhc_tcr}). The daily-mean field keeps a tight closed vortex because the storm
was still close to its landfall point during most of the day. The fastest winds (12.0\,m/s) are on its east side on the Atlantic coast near the Florida--Georgia border, where southerly flow over
the peninsula curves north-west.
(3) \textbf{Francine, D6.} The vorticity maximum ($15\times10^{-5}$\,s$^{-1}$) is at 30.23\,$^\circ$N, 88.31\,$^\circ$W on the
Mississippi coast, and the maximum speed (9.8\,m/s) is over the Louisiana delta. Inland, the field is dominated by a uniform
easterly flow over Mississippi, Alabama and Georgia. Francine is the weakest of the four in gridMET daily means.
(4) \textbf{Helene, D7.} The centre passed over the Gulf and therefore lies outside the land mask. The land shows only the
\emph{eastern half} of the vortex: 11--16\,m/s south-south-easterlies over the whole peninsula turning to south-easterlies
over Georgia and the Carolinas. The land vorticity maximum is at the south-west Florida coast (26.57\,$^\circ$N), which is a
shear line at the coast, not the eye. This is a key limitation of a land-only dataset for tropical cyclones: the
circulation centre must be inferred from the curvature of the flow.
(5) Daily averaging smears a moving storm along its track. Peak daily means of 10--16\,m/s are far below sustained
hurricane-force winds (at least 33\,m/s), so these fields show the storm-scale circulation, not the hazard.

\subsection{Streamlines vs.\ quiver}
Streamlines are integral curves of the vector field, tangent to $\mathbf{v}$ everywhere. Matplotlib's \texttt{streamplot}
integrates them on a regular grid. Because the map is projected, we regrid $(u,v)$ (rotated into the map frame) onto a
regular grid in projected metres (360$\times$224 for CONUS) and mask points outside the land mask, so streamlines stop at
the coast. Cartopy's built-in regridding (30$\times$density points) was too coarse. One rendering pitfall was found and
fixed: Matplotlib's streamline arrowheads were drawn with wrong orientation on a GeoAxes. We verified segment directions
numerically (0 of 74 sampled segments reversed) and redraw the heads as quiver arrows in map coordinates.

\figfull{Fig_22_streamlines_vs_quiver.png}{D7: (a) quiver of 20$\times$20 block means over the speed field; (b) streamlines
(density 2.4, field coarsened 3$\times$3) with line width and colour proportional to speed.}{fig:22}

\textbf{Inferences.}
(1) Streamlines show \emph{continuity} that quivers leave to the reader. In Fig.~\ref{fig:22}b one can follow a single
air stream from the Gulf coast of Texas north through Oklahoma, Kansas and Nebraska to the Dakotas: the Great Plains
southerly flow. One can also follow Helene's inflow from the Florida peninsula curving west and north-west across Georgia and
the Carolinas towards Tennessee.
(2) Streamlines expose \emph{shear and rotation}. Along about 95\,$^\circ$W a sharp shear zone separates the northward Plains
flow over Kansas and Nebraska from southward flow over Missouri and Arkansas: neighbouring streamlines run in opposite
directions. Small closed eddies appear
over the Great Basin and Four Corners region, and over northern Lake Michigan. In the quiver plot these features appear only as
inconsistent arrow directions.
(3) What streamlines hide: line density reflects the seeding algorithm, not speed, so magnitude must be coded separately
(colour and width here). Values at a given location cannot be read, and in the slow interior West (below 3\,m/s) the lines
wander and suggest more coherent flow than the 4\,km data support. Quivers remain better for local, quantitative
reading at known positions.

\subsection{Streamline seeding strategies}
Fig.~\ref{fig:23} compares three seeding strategies with the \emph{same number of seeds} (112) in the Helene zoom, using
unbroken streamlines integrated in both directions from each seed~\cite{turk1996,jobard1997}.
\begin{itemize}
\item \emph{Uniform grid}: a 15$\times$13 lattice of land points.
\item \emph{Random}: uniform on land cells, seed 7.
\item \emph{Feature-based}: 88 seeds drawn with probability $\propto v_s^4$ plus a 24-seed ring of radius 1.2$^\circ$
  around the maximum of smoothed cyclonic vorticity.
\end{itemize}
\figfull{Fig_23_streamline_seeding.png}{Seeding strategies on D7 (Helene zoom), 112 seeds each (red dots); streamline
colour = speed, grey background = speed. Coverage = share of land cells within about 25\,km of a streamline.}{fig:23}
\input{generated/tab_s3_seed}

\textbf{Inferences.}
(1) Uniform seeds give the most even coverage (98.9\,\% of land) and an unbiased overview. Every region gets lines, including
the slow flow over Alabama and Kentucky, but the high-speed core is covered last (92.6\,\%).
(2) Random seeds give similar coverage (98.0\,\%) with visible clumping, for example several seeds close together near the Mississippi River at about 35\,$^\circ$N
that duplicate one another while gaps open elsewhere, and results change with the seed. They are useful only to show that
conclusions do not depend on a lattice.
(3) Feature-based seeds concentrate on the flow that matters, covering 97.8\,\% of the top-10\,\% speed cells with 48\,\% fewer
line segments. The ring makes the strong south-south-easterly current over the peninsula and its curvature toward the Big
Bend explicit. The price is 77\,\% land coverage: slow regions are left empty, which is fine for an event narrative but
misleading as a general overview.

\subsection{Wind evolution over the eight dates}
\figfull{Fig_24_wind_evolution.png}{10\,m wind on the eight dates: speed (fixed 0--14\,m/s) and 24$\times$24 block-mean vectors
(about 96\,km). All eight frames of the GIF \fpath{Fig_24_wind_evolution.gif}, which uses 20$\times$20 blocks.}{fig:24}
\input{generated/tab_s3_dates}

\textbf{Inferences.}\footnote{Animation: \fpath{images/Fig_24_wind_evolution.gif} (8 frames, 1 frame/s, fixed speed colour bar).}
(1) Mean CONUS speed rises through the season from 2.8\,m/s (D1) to 3.9--4.0\,m/s (D6--D8) (Table~\ref{tab:s3dates}), as
the westerlies strengthen and migrate south in early autumn. The 99th percentile jumps to 10.1--10.3\,m/s only on the
Helene days.
(2) The share of cells with an eastward component falls from 62\,\% (D1) to 41\,\% (D5, D6), and the mean meridional
component is northward on seven dates. Late summer over CONUS is dominated by southerly return flow around the
Atlantic subtropical high, not by westerlies, which a pure ``prevailing westerlies'' expectation would get wrong.
(3) Event signatures: Beryl's cyclone over east Texas and Louisiana (D1); a separate cyclone over Minnesota and Wisconsin plus
Debby in Florida (D3); 5\,m/s mean southerly flow over the northern Plains (D5, 35--48\,$^\circ$N, 95--102\,$^\circ$W).
(4) The clearest signal is on D8. Helene's remnant is a large closed counter-clockwise gyre centred over western Kentucky and
Tennessee (vorticity maximum near 36.6\,$^\circ$N, 87.2\,$^\circ$W) with 10--14\,m/s winds over Georgia and the Carolinas.
Between landfall (about 10~miles west-southwest of Perry, about 30.0\,$^\circ$N 83.7\,$^\circ$W~\cite{nhc_tcr}) and D8 the vortex has moved about 800\,km inland, which explains the inland flooding rain in Fig.~\ref{fig:09}h.

\subsection{Observations and limitations (S3)}
Block averaging before subsampling, uniform arrows plus colour for pattern, proportional arrows for magnitude, streamlines
for continuity and topology, and feature seeding for event narratives together give a complete toolkit. Limitations: daily
means lose diurnal reversals such as sea breezes and mountain--valley winds, and hurricane peaks. The wind source is coarser
than the 4\,km grid. Land-only coverage cuts every hurricane in half at the coast.
